% ---------- Descriptive statistics (Krempl part 3) ----------
\begin{topic}{4 \; Descriptive statistics}{cStats}
\insight{\term{Descriptive} stats describe the sample as it is (descriptive analytics); \term{inferential} stats generalise to the population (predictive). Summaries lose information: always ask what a summary (KPI) can and cannot tell.}
\sub{Random variables \& distributions}
\insight{Discrete (die: $P(X{=}x_i)=1/6$, probability mass function) vs continuous (return of investment, density $f(x)$). Frequentist: how often it occurred; Bayesian: degree of belief.}
\formula{CDF}{F(x)=P(X\le x),\quad P(a\le X\le b)=F(b)-F(a),\quad P(X>a)=1-F(a)}{}
\insight{PDF $f(x)$ = density (probability of an infinitesimal interval around $x$), not a probability; for a normal it is highest at $\mu$.}
\ex{sales $\sim N(\mu{=}230\text{k},\sigma{=}20\text{k})$: $P(X>200\text{k})=1-\Phi(-1.5)=0.933$; $P(200\text{k}<X<260\text{k})=\Phi(1.5)-\Phi(-1.5)=0.866$; density at 230k $>$ at 240k.}
\begin{tabularx}{\linewidth}{@{}l*{8}{>{\centering\arraybackslash}X}@{}}
$z$ & 0 & 0.5 & 1 & 1.5 & 1.645 & 1.96 & 2 & 3\\
$\Phi(z)$ & .500 & .691 & .841 & .933 & .950 & .975 & .977 & .999\\
\end{tabularx}
\insight{$\Phi(-z)=1-\Phi(z)$.}
\sub{Central tendency (``centre of gravity'')}
\insight{\term{Mode} most frequent (nominal+); a distinct mode in real-valued data can reveal spurious values (0 = missing); meaningless for continuous unique values. \term{Median} 50\% below/above (ordinal+), robust.}
\formula{Arithmetic mean}{\xbar=\frac{1}{n}\sum_{i=1}^n x_i}{interval+; \term{susceptible to outliers} (Gaussian+outlier: median 10, mean 35).}
\formula{Geometric mean}{\Big(\prod_{i=1}^n x_i\Big)^{1/n}}{ratio scale only.}
\sub{Dispersion}
\formula{Range}{\max x-\min x}{}
\formula{IQR}{Q_3-Q_1}{$Q_1$: 25\% smaller; $Q_3$: 75\% smaller; robust (ignores outliers) unlike variance; related: percentiles.}
\formula{Variance}{s^2=\frac{\sum_{i=1}^n(x_i-\xbar)^2}{n-1}}{$n-1$: unbiased sample variance (dividing by $n$: 2nd moment around mean); $s=\sqrt{s^2}$.}
\formula{Mean abs.\ deviation}{\text{MAD}=\frac{1}{n}\sum|x_i-\xbar|}{penalises large deviations less than $s$.}
\sub{Shape}
\formulaD{Skewness (symmetry)}{\text{skew}=\frac{\frac1n\sum(x_i-\xbar)^3}{\big(\frac1n\sum(x_i-\xbar)^2\big)^{3/2}}\cdot\frac{n^2}{(n-1)(n-2)}}{$>0$: \term{tail on the right}, mass ``swayed left'', usually mean $>$ median. $<0$: \term{tail on the left}, usually mean $<$ median. $\approx0$: symmetric. Exceptions exist (discrete data, Weibull): the cube makes the \emph{longer tail} decide, not the mean--median order. Book p.106 is wrong.}
\formula{Kurtosis}{\frac{\sum(x_i-\xbar)^4}{n\,s^4}\ (\text{normal}=3);\ \text{excess}=\text{kurt}-3}{peakedness vs normal: excess $=0$ mesokurtic (normal), $>0$ leptokurtic (tall/sharp peak), $<0$ platykurtic (flat). E.g.\ uniform 1.8, bimodal 1.44, peaked 17.8, with outlier 98.}
\formula{Moments}{\text{moment}_p=E[(X-\xbar)^p]}{$p=2$ variance, 3 skewness, 4 kurtosis (standardised).}
\ex{$X_1=\{-2,-1,-.5,0,.25,.25,1,2\}$ symmetric; $X_2=\{0,.1,.3,.7,1.5,3,7\}$ positive skew, mean $>$ median; $X_3=\{-2,-1.25,-1,-.8,-.5,0,0,0\}$ bimodal (mode 0).}
\sub{Information content}
\formula{Entropy}{H(X)=-\sum_{x}P(x)\log P(x),\ \ 0\log0=0}{higher = more uncertain / informative outcome; $\log_2$ $\to$ bits: fair coin = 1 bit (max of a binary var.\ at $p=0.5$); uniform data highest.}
\formula{KL divergence}{D(p\|q)=\sum_x p(x)\log\frac{p(x)}{q(x)}}{divergence of approximation $q$ from true $p$; $\ge0$, $=0$ iff $p=q$; \term{not symmetric}, no triangle inequality $\Rightarrow$ not a true distance; $=\infty$ if $p(x)>0$ and $q(x)=0$; $0\log\frac0q=0$.}
\ex{shares $P$=(.4,.3,.3), NL (.4,.35,.25), IT (.5,.3,.2), ln: $D(P\|NL)=0.0085$, $D(P\|IT)=0.032$ $\Rightarrow$ NL more similar to world-wide.}
\sub{Two variables}
\formula{Covariance}{\text{cov}(x,y)=\frac{\sum(x_i-\xbar)(y_i-\bar y)}{n-1}}{sign = direction of linear relation; scale-dependent.}
\formula{Pearson}{r=\frac{\text{cov}(x,y)}{s_x s_y}\in[-1,1]}{linear relation only: $\pm1$ perfect linear, $0$ no \emph{linear} relation (variables can still be dependent!).}
\insight{\term{Spearman}: rank the values, compute Pearson on the ranks $\Rightarrow$ measures monotone relation (non-linear monotone: Spearman 1.00 while Pearson 0.30). Correlation $\ne$ causality.}
\formula{Independence}{P(A,B)=P(A)\,P(B)}{knowing $A$ does not help predict $B$.}
\formula{Cond.\ independence}{P(A,B|C)=P(A|C)P(B|C)\Leftrightarrow P(A|B,C)=P(A|C)}{does \emph{not} imply unconditional independence (butter \& milk dependent, independent given lactose tolerance).}
\sub{Linear regression}
\formula{Simple}{y=\beta_0+\beta_1x+\epsilon}{$\beta_0$ intercept, $\beta_1$ slope (change in $y$ per unit $x$).}
\formula{Multiple}{y=\beta_0+\beta_1x_1+\dots+\beta_px_p}{simple: 1 explanatory var; multiple: several; multivariate: several dependent vars.}
\insight{Fit by \term{OLS}: minimise sum of squared errors $\sum(y_i-\hat y_i)^2$. Purpose: descriptive (influence of $X$ on $Y$) or predictive.}
\insight{Assumptions (violation $\to$ remedy): \term{linearity}; \term{independence} of observations (same person twice $\to$ keep one per person); \term{normal residuals} (skewed $\to$ non-linear transform of $Y$); \term{homoscedasticity} = constant residual variance (funnel $\to$ transform $Y$); \term{no multicollinearity} (age \& birth year $\to$ unreliable coefficients $\to$ combine/drop).}
\sub{Point vs interval estimates, testing}
\insight{All above are \term{point estimates}: one value, no info on how well it describes the data (sample size!). Inferential: confidence intervals \& tests. One-sample t-test: $H_1:\mu>40$, $H_0:\mu\le40$; $\alpha$ = P(wrongly rejecting $H_0$); reject $H_0$ only if p-value $<\alpha$ (p $=0.5$ $\Rightarrow$ $H_1$ not confirmed).}
\end{topic}
